% Part: computability
% Chapter: recursive-functions
% Section: examples

\documentclass[../../../include/open-logic-section]{subfiles}

\begin{document}

\olfileid{cmp}{rec}{exa}
\olsection{أمثلة للدوال العودية البدائية}

تقدمت لنا أمثلة للدوال العودية البدائية، منها الجمع
والضرب~$\Add$ و$\Mult$. ومنها دالة الهوية $\fn{id}(x)=x$،
فهي عين~$\Proj{1}{0}$. وكذلك الدوال الثابتة
$\fn{const}_n(x)=n$؛ فإنها تستخرج من $\Zero$ و$\Succ$
بالتركيب مرة بعد مرة. ونحتاج إلى هذه الدوال حين نورد الثوابت في التعريف
العودي البدائي. فمن طلب، مثلًا، تعريف $f(x)=2\cdot x$، استخرجها
بتركيب $\fn{const}_2(x)$ مع الضرب، فيضع:
$f(x)=\Mult(\fn{const}_2(x),\Proj{1}{0}(x))$. وسنجري على هذه الطريقة في
ما يأتي.

\begin{prop}
  دالة الأس $\fn{exp}(x,y)=x^y$ عودية بدائية.
\end{prop}

\begin{proof}
  تعريف $\fn{exp}$ بالعودية البدائية على الوجه الآتي:
  \begin{align*}
    \fn{exp}(x, 0) & = 1\\
    \fn{exp}(x, y+1) & = \Mult(x, \fn{exp}(x,y)).
    \intertext{وهذا، بالدقة الصارمة، ليس تعريفًا عوديًا من دوال عودية
      بدائية. أما رسميًا فلدينا:}
    \fn{exp}(x, 0) & = f(x)\\
    \fn{exp}(x, y+1) & = g(x, y, \fn{exp}(x,y)).
    \intertext{حيث}
    f(x) & = \Succ(\Zero(x)) = 1\\
    g(x, y, z) & = \Mult(\Proj{3}{0}(x, y, z), \Proj{3}{2}(x, y, z)) = x \cdot z
  \end{align*}
  فقد حصلت $f$ و$g$ بالتركيب من دوال عودية بدائية، وهو المطلوب.
\end{proof}

\begin{prop}
  دالة السابق $\fn{pred}(y)$ المعرفة بـ
  \[
  \fn{pred}(y) = \begin{cases}
    0 & \text{إذا كان $y=0$}\\
    y-1 & \text{في غير ذلك}
  \end{cases}
  \]
  عودية بدائية.
\end{prop}

\begin{proof}
 أصل الحجة أن
 \begin{align*}
   \fn{pred}(0) & = 0 \text{، وأن}\\
   \fn{pred}(y+1) & = y.
 \end{align*}
 وهذا قريب من التعريف بالعودية البدائية، غير أنه لا يوافق
 نمطه الرسمي؛ إذ لا بد في ذلك النمط من حجة إضافية
 واحدة على الأقل~$x$. ثم إن القيمة $\fn{pred}(y)$ لا تدخل
 في تعريف $\fn{pred}(y+1)$. ويُستدرك الأمر بتعريف
 $\fn{pred}'(x,y)$ أولًا هكذا:
 \begin{align*}
   \fn{pred}'(x, 0) & = \Zero(x) = 0,\\
   \fn{pred}'(x, y+1) & = \Proj{3}{1}(x, y, \fn{pred'}(x, y)) = y.
 \end{align*}
ثم نستخرج $\fn{pred}$ منها بالتركيب، فنضع مثلًا:
$\fn{pred}(x)=\fn{pred}'(\Zero(x),\Proj{1}{0}(x))$.
\end{proof}

\begin{prop}
  دالة المضروب $\fn{fac}(x)=\fact{x}=1\cdot2\cdot3
  \cdot\dots\cdot x$ عودية بدائية.
\end{prop}

\begin{proof}
  يظهر تعريفها بالعودية البدائية من المعادلتين
  \begin{align*}
    \fn{fac}(0) &= 1\\
    \fn{fac}(y+1) & = \fn{fac}(y) \cdot (y+1).
    \intertext{ويلزمنا رسميًا أن نعرف أولًا دالة ثنائية الحجج $h$:}
    h(x, 0) & = \fn{const}_1(x)\\
    h(x, y+1) & = g(x, y, h(x, y))
    \intertext{حيث $g(x,y,z)=\Mult(\Proj{3}{2}(x,y,z),
      \Succ(\Proj{3}{1}(x,y,z)))$، ثم نجعل}
    \fn{fac}(y) & = h(\Proj{1}{0}(y), \Proj{1}{0}(y)) = h(y,y).
  \end{align*}
  وقد اتضح وجه الرد إلى الصورة الرسمية؛ فلن نلتزم في ما يأتي بذكر كل تفصيل
  من تفاصيل التركيب والعودية البدائية.
\end{proof}

\begin{prop}
  الطرح المبتور، $x\tsub y$، المعرف بـ
  \[
  x \tsub y = \begin{cases}
    0 & \text{إذا كان $x<y$}\\
    x-y & \text{في غير ذلك}
  \end{cases}
  \]
  عودي بدائي.
\end{prop}

\begin{proof}
  لدينا:
  \begin{align*}
    x \tsub 0 & = x\\
    x \tsub (y+1) & = \fn{pred}(x \tsub y)
  \end{align*}
\end{proof}

\begin{prop}
 المسافة بين $x$ و$y$، أي $\left|x-y\right|$، عودية بدائية.
\end{prop}

\begin{proof}
  لما كان $\left|x-y\right|=(x\tsub y)+(y\tsub x)$، صح تعريف
  المسافة بتركيب $+$ و$\tsub$، وقد ثبتت عوديتهما البدائية.
\end{proof}

\begin{prop}
  أكبر $x$ و$y$، أي $\fn{max}(x,y)$، عودي بدائي.
\end{prop}

\begin{proof}
  نستخرج $\fn{max}(x,y)$ من $+$ و$\tsub$ بالتركيب الآتي:
  \[
  \fn{max}(x,y) \defis x + (y \tsub x).
  \]
  وبيانه: إن كان $x$ أكبرهما، أي $x\ge y$، لزم $y\tsub x=0$، فكان
  $x+(y\tsub x)=x+0=x$. وإن كان $y$ أكبرهما، لزم
  $y\tsub x=y-x$، ومن ثم $x+(y\tsub x)=x+(y-x)=y$.
\end{proof}

\begin{prop}
  \ollabel{prop:min-pr}
  أصغر $x$ و$y$، أي $\fn{min}(x,y)$، عودي بدائي.
\end{prop}

\begin{proof}
  تمرين.
\end{proof}

\begin{prob}
  أقم البرهان على \olref[cmp][rec][exa]{prop:min-pr}.
\end{prob}

\begin{prob}
بين أن
\[
f(x, y) =
2^{(2^{\iddots^{2^{x}}})}\raisebox{1ex}{\bigg\rbrace}
\raisebox{1ex}{\text{$y$ نسخ من العدد~$2$}}
\]
دالة عودية بدائية.
\end{prob}

\begin{prob}
أثبت أن القسمة الصحيحة $d(x,y)=\lfloor x/y\rfloor$، وهي القسمة التي نطرح فيها
كل ما بعد الفاصلة العشرية، عودية بدائية. واصطلاحنا عند $y=0$ أن
$d(x,y)=0$. واكتب لـ~$d$ تعريفًا صريحًا بالعودية البدائية
والتركيب.
\end{prob}

\begin{prop}
لا تخرج العمليتان الآتيتان عن مجموعة الدوال العودية البدائية:
\begin{enumerate}
\item المجاميع المنتهية: إذا كانت $f(\vec x,z)$ عودية بدائية، كانت
الدالة الآتية كذلك:
\[
g(\vec x, y) \defis \sum_{z = 0}^y f(\vec x, z).
\]
\item الجداءات المنتهية: إذا كانت $f(\vec x,z)$ عودية بدائية، كانت
الدالة الآتية كذلك:
\[
h(\vec x, y) \defis \prod_{z = 0}^y f(\vec x, z).
\]
\end{enumerate}
\end{prop}

\begin{proof}
أما المجاميع المنتهية، فتعريفها العودي، مثلًا، بالمعادلتين
\begin{align*}
  g(\vec x, 0) & = f(\vec x, 0)\\
  g(\vec x, y+1) & = g(\vec x, y) + f(\vec x, y+1).
\end{align*}
\end{proof}


\end{document}
